% Cover letter -- InkCloser, submission to Computer Vision and Image Understanding
% Compile: pdflatex cover_letter.tex
\documentclass[11pt]{article}
\usepackage[a4paper,margin=2.3cm]{geometry}
\usepackage[T1]{fontenc}
\usepackage[utf8]{inputenc}
\usepackage{parskip}
\usepackage[hidelinks]{hyperref}
\pagestyle{empty}

\begin{document}

\begin{flushright}
Ruben Tous\\
Computer Architecture Department\\
Universitat Polit\`ecnica de Catalunya\\
Barcelona, Spain\\
\href{mailto:ruben.tous@upc.edu}{ruben.tous@upc.edu}\\[0.8em]
\today
\end{flushright}

Editor-in-Chief\\
\emph{Computer Vision and Image Understanding}

Dear Editor,

I am pleased to submit the manuscript entitled \textbf{``InkCloser: Line-Art Completion
for Paint-Bucket Colorization''} for consideration as a research article in
\emph{Computer Vision and Image Understanding}.

Flood-fill (paint-bucket) coloring, the standard way of coloring line art in
illustration and 2D animation, only works when the drawing forms closed regions; a
single gap in a stroke makes the fill leak into the neighboring region. The manuscript
addresses the restoration of such gaps as a \emph{mask-free} problem: unlike
conventional inpainting, the location of the missing ink is unknown at inference time,
so the network must find the gaps, close them, and leave intentional stroke endings and
the rest of the drawing untouched.

The main contributions are:
\begin{itemize}\setlength{\itemsep}{1pt}
\item a U-Net with gated convolutions, a dilated context stack, and transformer
  self-attention at two resolutions, designed to locate and bridge gaps without a mask;
\item a loss that combines hole-weighted cross-entropy with global and gap-local Dice
  terms;
\item evaluation metrics that measure gap closure and simulate the paint-bucket tool
  (bucket accuracy and fill leaks), rather than relying on pixel accuracy;
\item a public dataset of 55\,000 corrupted and clean line-art pairs, together with the
  code and pretrained weights.
\end{itemize}

On a held-out test set of 1\,000 drawings, the model closes 70.5\% of the gaps, raises
the fraction of regions that a single bucket click fills correctly from 64.1\% to
94.7\%, reduces fill leaks per image from 15.5 to 1.5, and makes 45.0\% of the drawings
entirely leak-free. Morphological closing and the joint gap detection and inpainting
network of Sasaki et al.\ (CVPR 2017), even when retrained on the same data, do not
improve on the bucket accuracy of the unrestored input.

I believe the work fits the scope of \emph{Computer Vision and Image Understanding}, as
it addresses image restoration and structural understanding of line drawings, with an
evaluation protocol tied to a concrete downstream task, and it is relevant to the
growing body of work on learning-based tools for illustration and animation.

I confirm that this manuscript is original, has not been published previously, and is
not under consideration for publication elsewhere. I am the sole author and approve the
submitted version. I have no competing interests to declare. The code, pretrained
weights, and dataset are publicly available at
\url{https://github.com/rtous/inkcloser}.
% Optional, if the PDF remains in the repository:
% A preprint of the manuscript is available in the same repository.

Thank you for your time and consideration.

Sincerely,\\[1.2em]
Ruben Tous\\
Universitat Polit\`ecnica de Catalunya

\end{document}
